% Options for packages loaded elsewhere
% Options for packages loaded elsewhere
\PassOptionsToPackage{unicode}{hyperref}
\PassOptionsToPackage{hyphens}{url}
\PassOptionsToPackage{dvipsnames,svgnames,x11names}{xcolor}
%
\documentclass[
  letterpaper,
  DIV=11,
  numbers=noendperiod]{scrartcl}
\usepackage{xcolor}
\usepackage{amsmath,amssymb}
\setcounter{secnumdepth}{-\maxdimen} % remove section numbering
\usepackage{iftex}
\ifPDFTeX
  \usepackage[T1]{fontenc}
  \usepackage[utf8]{inputenc}
  \usepackage{textcomp} % provide euro and other symbols
\else % if luatex or xetex
  \usepackage{unicode-math} % this also loads fontspec
  \defaultfontfeatures{Scale=MatchLowercase}
  \defaultfontfeatures[\rmfamily]{Ligatures=TeX,Scale=1}
\fi
\usepackage{lmodern}
\ifPDFTeX\else
  % xetex/luatex font selection
\fi
% Use upquote if available, for straight quotes in verbatim environments
\IfFileExists{upquote.sty}{\usepackage{upquote}}{}
\IfFileExists{microtype.sty}{% use microtype if available
  \usepackage[]{microtype}
  \UseMicrotypeSet[protrusion]{basicmath} % disable protrusion for tt fonts
}{}
\makeatletter
\@ifundefined{KOMAClassName}{% if non-KOMA class
  \IfFileExists{parskip.sty}{%
    \usepackage{parskip}
  }{% else
    \setlength{\parindent}{0pt}
    \setlength{\parskip}{6pt plus 2pt minus 1pt}}
}{% if KOMA class
  \KOMAoptions{parskip=half}}
\makeatother
% Make \paragraph and \subparagraph free-standing
\makeatletter
\ifx\paragraph\undefined\else
  \let\oldparagraph\paragraph
  \renewcommand{\paragraph}{
    \@ifstar
      \xxxParagraphStar
      \xxxParagraphNoStar
  }
  \newcommand{\xxxParagraphStar}[1]{\oldparagraph*{#1}\mbox{}}
  \newcommand{\xxxParagraphNoStar}[1]{\oldparagraph{#1}\mbox{}}
\fi
\ifx\subparagraph\undefined\else
  \let\oldsubparagraph\subparagraph
  \renewcommand{\subparagraph}{
    \@ifstar
      \xxxSubParagraphStar
      \xxxSubParagraphNoStar
  }
  \newcommand{\xxxSubParagraphStar}[1]{\oldsubparagraph*{#1}\mbox{}}
  \newcommand{\xxxSubParagraphNoStar}[1]{\oldsubparagraph{#1}\mbox{}}
\fi
\makeatother


\usepackage{longtable,booktabs,array}
\newcounter{none} % for unnumbered tables
\usepackage{calc} % for calculating minipage widths
% Correct order of tables after \paragraph or \subparagraph
\usepackage{etoolbox}
\makeatletter
\patchcmd\longtable{\par}{\if@noskipsec\mbox{}\fi\par}{}{}
\makeatother
% Allow footnotes in longtable head/foot
\IfFileExists{footnotehyper.sty}{\usepackage{footnotehyper}}{\usepackage{footnote}}
\makesavenoteenv{longtable}
\usepackage{graphicx}
\makeatletter
\newsavebox\pandoc@box
\newcommand*\pandocbounded[1]{% scales image to fit in text height/width
  \sbox\pandoc@box{#1}%
  \Gscale@div\@tempa{\textheight}{\dimexpr\ht\pandoc@box+\dp\pandoc@box\relax}%
  \Gscale@div\@tempb{\linewidth}{\wd\pandoc@box}%
  \ifdim\@tempb\p@<\@tempa\p@\let\@tempa\@tempb\fi% select the smaller of both
  \ifdim\@tempa\p@<\p@\scalebox{\@tempa}{\usebox\pandoc@box}%
  \else\usebox{\pandoc@box}%
  \fi%
}
% Set default figure placement to htbp
\def\fps@figure{htbp}
\makeatother





\setlength{\emergencystretch}{3em} % prevent overfull lines

\providecommand{\tightlist}{%
  \setlength{\itemsep}{0pt}\setlength{\parskip}{0pt}}



 


\KOMAoption{captions}{tableheading}
\makeatletter
\@ifpackageloaded{caption}{}{\usepackage{caption}}
\AtBeginDocument{%
\ifdefined\contentsname
  \renewcommand*\contentsname{Table of contents}
\else
  \newcommand\contentsname{Table of contents}
\fi
\ifdefined\listfigurename
  \renewcommand*\listfigurename{List of Figures}
\else
  \newcommand\listfigurename{List of Figures}
\fi
\ifdefined\listtablename
  \renewcommand*\listtablename{List of Tables}
\else
  \newcommand\listtablename{List of Tables}
\fi
\ifdefined\figurename
  \renewcommand*\figurename{Figure}
\else
  \newcommand\figurename{Figure}
\fi
\ifdefined\tablename
  \renewcommand*\tablename{Table}
\else
  \newcommand\tablename{Table}
\fi
}
\@ifpackageloaded{float}{}{\usepackage{float}}
\floatstyle{ruled}
\@ifundefined{c@chapter}{\newfloat{codelisting}{h}{lop}}{\newfloat{codelisting}{h}{lop}[chapter]}
\floatname{codelisting}{Listing}
\newcommand*\listoflistings{\listof{codelisting}{List of Listings}}
\makeatother
\makeatletter
\makeatother
\makeatletter
\@ifpackageloaded{caption}{}{\usepackage{caption}}
\@ifpackageloaded{subcaption}{}{\usepackage{subcaption}}
\makeatother
\usepackage{bookmark}
\IfFileExists{xurl.sty}{\usepackage{xurl}}{} % add URL line breaks if available
\urlstyle{same}
\hypersetup{
  pdftitle={Midterm 2},
  pdfauthor={Arun, David, and Daniel},
  colorlinks=true,
  linkcolor={blue},
  filecolor={Maroon},
  citecolor={Blue},
  urlcolor={Blue},
  pdfcreator={LaTeX via pandoc}}


\title{Midterm 2}
\usepackage{etoolbox}
\makeatletter
\providecommand{\subtitle}[1]{% add subtitle to \maketitle
  \apptocmd{\@title}{\par {\large #1 \par}}{}{}
}
\makeatother
\subtitle{Taxation, Redistribution, and the Macroeconomy}
\author{Arun, David, and Daniel}
\date{November 22, 2026}
\begin{document}
\maketitle


\section{Question 1: David}\label{question-1-david}

\subsection{Part A (Mirrlees)}\label{part-a-mirrlees}

\begin{enumerate}
\def\labelenumi{\alph{enumi}.}
\tightlist
\item
  The planner solves:
\end{enumerate}

\[ 
\max_{c_L,c_H,y_L,y_H}\ \tfrac12\Big[u(c_L)-v\big(\tfrac{y_L}{w_L}\big)\Big]+\tfrac12\Big[u(c_H)-v\big(\tfrac{y_H}{w_H}\big)\Big]\quad\text{s.t.}\quad \tfrac12 c_L+\tfrac12 c_H \le \tfrac12 y_L+\tfrac12 y_H .
\]

FOCS:

\[
\begin{aligned}
[c_l]:&\quad \tfrac12 u'(c_l)-\tfrac12\lambda=0 &&\rightarrow\quad u'(c_l)=\lambda \\
[y_l]:&\quad -\tfrac12\,\frac{v'(y_l/w_l)}{w_l}+\tfrac12\lambda=0 &&\rightarrow\quad \frac{v'(y_l/w_l)}{w_l} = \lambda \\
[c_h]:&\quad \tfrac12 u'(c_h)-\tfrac12\lambda=0 &&\rightarrow\quad u'(c_h)=\lambda \\
[y_h]:&\quad -\tfrac12\,\frac{v'(y_h/w_h)}{w_h}+\tfrac12\lambda=0 &&\rightarrow\quad \frac{v'(y_h/w_h)}{w_h} = \lambda \\
[\lambda]:&\quad y_l + y_h - c_l - c_h = 0
\end{aligned}
\]

Like in class, after we combine these FOCS, we get the nice equations
that MRS=MRT.

\[
\begin{aligned}
\lambda = \lambda \quad &\rightarrow \quad u'(c_L)=u'(c_H) \\
\lambda = \lambda \quad &\rightarrow \quad \frac{v'(n_l)}{w_l} = \frac{v'(n_h)}{w_h}
\end{aligned}
\]

Both types consume the same as per the first equation. High type has a
higher wage given their higher productivity, and higher labor supply,
thus their income is greater than the low-type. In this setting, the
more productive type works more and earns more but consumes the same, so
the high type is a net taxpayer and the low type a net recipient.

\begin{enumerate}
\def\labelenumi{\alph{enumi}.}
\setcounter{enumi}{1}
\item
  If you're the high-type, and you're more productive, and yet consuming
  the same of the low-type, then the high-type would choose the low-type
  bundle. The low-type will also choose the low-type bundle. This is a
  problem because if both agents pick \((c_L,y_L)\), then earnings fall,
  and the resource constraint is violated. The first best is not
  incentive compatible, which is why we need incentive constraints.
\item
  I'll put the planners max. problem in a nicer format, as the \(1/2\)
  eventually goes away.
\end{enumerate}

\[
\begin{aligned}
\max_{c_l,c_h, n_l, n_h}\ & \Big[u(c_L)-v\big(\tfrac{y_L}{w_L}\big)\Big]+\Big[u(c_H)-v\big(\tfrac{y_H}{w_H}\big)\Big]\\
\text{s.t.}\ & c_L+c_H= y_L+y_H \\
& v_H(c_H, y_H)\ = \ v_H(c_L,y_L) &&\text{(IC}_H), \\
& v_L(c_L, y_L)\ > \ v_L(c_H,y_H) &&\text{(IC}_L)
\end{aligned}
\]

The high-type IC says the high-type must prefer their bundle to
mimicking the low-type. The low-type IC says the same for the low-type.
These only rule out allocations, so the planner can't implement the
first best as in part b where the high-type would mimic.

At the separating optimum, the high-type IC binds, since the planner
redistributes until the high-type is just indifferent to mimicking. The
low-type IC is slack, since the low-type never wants the high-type's
bundle (more hours at a lower wage). The resource constraint binds,
since nothing is wasted.

\begin{enumerate}
\def\labelenumi{\alph{enumi}.}
\setcounter{enumi}{3}
\tightlist
\item
  In class we've stated the following conclusion:
\end{enumerate}

\begin{quote}
\begin{enumerate}
\def\labelenumi{\arabic{enumi}.}
\tightlist
\item
  At an optimum, we can't take all resources from being productive.
  \(V_H > V_L\). The planner can't equalize utilities. Mimicking takes
  fewer hours, so that utility is always more than the low-type gets.
\end{enumerate}
\end{quote}

\begin{quote}
\begin{enumerate}
\def\labelenumi{\arabic{enumi}.}
\setcounter{enumi}{1}
\tightlist
\item
  At a planner efficient condition, we want the high-type to work as
  much as possible, so we don't distrubt their labor supply. We don't
  distort the high-type's labor supply, \(T'(y_H)=0\). The high-type
  works the efficient amount (MRS=MRT). Taxing \(y_H\) at the margin
  wouldn't stop mimicking, so it would only cost efficiency. This
  doesn't mean the high-type pays no taxes as their average tax is still
  positive.
\end{enumerate}
\end{quote}

\begin{quote}
\begin{enumerate}
\def\labelenumi{\arabic{enumi}.}
\setcounter{enumi}{2}
\tightlist
\item
  We can have a marginal tax (T'(y\_L)\textgreater0) to preserve
  incentives, so we can't make the low-type attractice so that the
  high-type does not end up mimicking the low-type. This relaxes the
  high-type IC and lets the planner redistribute more.
\end{enumerate}
\end{quote}

\subsection{Part B (Ramsey)}\label{part-b-ramsey}

\breakpage

\section{Q1: Arun}\label{q1-arun}

\subsection{Part A: Mirrlees Taxation}\label{part-a-mirrlees-taxation}

\subsubsection{a:}\label{a}

\[
\begin{aligned}
\max_{c_L,c_H,y_L,y_H} \quad
& u(c_L)+u(c_H)
-v\left(\frac{y_L}{w_L}\right)
-v\left(\frac{y_H}{w_H}\right) \\
\text{s.t.} \quad
& c_L+c_H=y_L+y_H.
\end{aligned}
\]

\[
\begin{aligned}
\mathcal{L}
= {} & u(c_L) + u(c_H)
- v\left(\frac{y_L}{w_L}\right)
- v\left(\frac{y_H}{w_H}\right) + \lambda(y_L+y_H-c_L-c_H).
\end{aligned}
\]

FOCs:

\[
\begin{aligned}
\frac{\partial\mathcal{L}}{\partial c_L}
&=u'(c_L)-\lambda=0,\\
\frac{\partial\mathcal{L}}{\partial c_H}
&=u'(c_H)-\lambda=0,\\
\frac{\partial\mathcal{L}}{\partial y_L}
&=-\frac{1}{w_L}v'\left(\frac{y_L}{w_L}\right)+\lambda=0,\\
\frac{\partial\mathcal{L}}{\partial y_H}
&=-\frac{1}{w_H}v'\left(\frac{y_H}{w_H}\right)+\lambda=0,\\
\frac{\partial\mathcal{L}}{\partial\lambda}
&=y_L+y_H-c_L-c_H=0.
\end{aligned}
\]

Combining the consumption FOCs gives:

\[
u'(c_L)=u'(c_H)=\lambda.
\]

Since \(u''<0\), marginal utility is decreasing in consumption:

\[
c_L=c_H=\bar{c}.
\]

Combining the earnings FOCs gives:

\[
\frac{v'(y_L/w_L)}{w_L}
=
\frac{v'(y_H/w_H)}{w_H}
=\lambda.
\]

Since \(n_i=y_i/w_i\), we can rewrite as:

\[
v'(n_L)=\lambda w_L,
\qquad
v'(n_H)=\lambda w_H.
\]

Since \(w_H>w_L\) and \(\lambda>0\),

\[
v'(n_H)>v'(n_L).
\]

Because \(v''>0\), marginal disutility of labor is strictly increasing,
so \(n_H>n_L\). Therefore, we can say that:

\[
y_H>y_L
\]

We can combine the consumption and earnings FOCs to get:

\[
\frac{v'(y_i/w_i)}{w_i}=u'(c_i).
\\\\
\frac{v'(n_i)}{u'(c_i)}=w_i.
\]

Via this, we have that the marginal rate of substitution equals the
marginal rate of transformation, so the first-best efficiency condition
is satisfied.

Since \(c_H=c_L=\bar{c}\), net taxes for each type are:

\[
t_i=y_i^*-\bar{c}.
\] Since \(y_H^*>y_L^*\), we have \(t_H>0\) and \(t_L<0\). So, there is
redistribution from the high-productivity to the low-productivity
workers.

\subsubsection{b:}\label{b}

For the low-productivity worker, since \(y_h>y_L\), choosing the
high-type bundle would require more labor without any additional
consumption, so they choose the low-type bundle. The utility functions
are identical, so the high-productivity worker would do the same.

\subsubsection{c:}\label{c}

\[
\begin{aligned}
\max_{c_L,c_H,y_L,y_H}\quad
&u(c_L)+u(c_H)
-v(y_L/w_L)-v(y_H/w_H)\\
\text{s.t.}\quad
&c_L+c_H=y_L+y_H,\\
\quad
&u(c_H)-v(y_H/w_H)
= u(c_L)-v(y_L/w_H), \quad \text{(IC-H)} \\
\quad
&u(c_L)-v(y_L/w_L)
> u(c_H)-v(y_H/w_L). \quad \text{(IC-L)}
\end{aligned}
\]

These conditions require the workers to prefer their intended bundles
(the high-productivity worker prefers the high-type bundle, and vice
versa). These constraints restrict the feasible allocations to those
under which neither type wants the other's bundle. The planner cannot
implement the first-best allocation because both types would choose the
low-type bundle.

At the separating optimum, the high-productivity constraint binds, while
the low-productivity constraint is slack. The former holds with
equality, while the latter holds with strict inequality.

\subsubsection{d:}\label{d}

The three results are:

\begin{enumerate}
\def\labelenumi{\arabic{enumi}.}
\tightlist
\item
  Tax levels implement the optimal allocation and redistribute. At the
  redistributive optimum, the high type pays net taxes and the low type
  receives a transfer. This follows from the budget constraint
  \(c_i=y_i-T(y_i)\).
\end{enumerate}

\[
T(y_i^*)=y_i^*-c_i^*.
\]

\begin{enumerate}
\def\labelenumi{\arabic{enumi}.}
\setcounter{enumi}{1}
\tightlist
\item
  No distortion at the top. The high type faces a zero marginal income
  tax rate, so additional earnings are not taxed at the margin. This
  follows from the optimality conditions. Distorting the high-type labor
  does not relax the high-type incentive constraint.
\end{enumerate}

\[
\boxed{T'(y_H^*)=0.}
\]

\begin{enumerate}
\def\labelenumi{\arabic{enumi}.}
\setcounter{enumi}{2}
\tightlist
\item
  Positive marginal taxation at the bottom. This happens because the
  high-type incentive constraint binds. Reducing the low type's earnings
  makes the bundle less attractive to high-productivity workers who
  might take the low type. This allows for more redistribution.
\end{enumerate}

\[
\boxed{0<T'(y_L^*)<1.}
\]

\subsection{Part B: Ramsey taxation}\label{part-b-ramsey-taxation}

\subsubsection{e:}\label{e}

\[
\begin{aligned}
\max_{c_i,n_i}\quad
&u(c_i)-v(n_i)\\
\text{s.t.}\quad
&c_i=w_i n_i-T(w_i n_i;\theta).
\end{aligned}
\]

The FOC for \(n_i\) is:

\[
u'(c_i)w_i[1-T_y(y_i;\theta)]-v'(n_i)=0.
\]

\[
\frac{v'(n_i)}{u'(c_i)}
=w_i[1-T_y(y_i;\theta)].
\]

When we take \(\theta=(t,B)\) and \(T(y;\theta)=ty-B\), the household
problem becomes:

\[
\begin{aligned}
\max_{c_i,n_i}\quad
&u(c_i)-v(n_i)\\
\text{s.t.}\quad
&c_i=(1-t)w_i n_i+B.
\end{aligned}
\]

The FOC under the above method simplifies to:

\[
\frac{v'(n_i)}{u'(c_i)}=w_i(1-t).
\]

An additional hour of work will increase earnings by \(w_i\), but sets
after-tax consumption to \(w_i(1-t)\).

So, the tax schedule will reduce the utility gotten from an additional
hour of work, and create a wedge between the MRS and productivity. The
first-best condition is not met unless the marginal tax rate is zero.

\subsection{f:}\label{f}

For the tax schedule \(T(y;\theta)\) consists of allocations
\((c_i,n_i,y_i)\) for each type \(i\in\{L,H\}\) such that:

\begin{enumerate}
\def\labelenumi{\arabic{enumi}.}
\tightlist
\item
  Households optimize. They do so according to:
\end{enumerate}

\[
\begin{aligned}
\max_{c_i,n_i}\quad
&u(c_i)-v(n_i)\\
\text{s.t.}\quad
&c_i=w_i n_i-T(w_i n_i;\theta).
\end{aligned}
\]

\begin{enumerate}
\def\labelenumi{\arabic{enumi}.}
\setcounter{enumi}{1}
\tightlist
\item
  The government budget balances. As the government does not make
  purchases, the aggregate net tax revenue must equal zero:
\end{enumerate}

\[
T(y_L;\theta)+T(y_H;\theta)=0.
\]

\begin{enumerate}
\def\labelenumi{\arabic{enumi}.}
\setcounter{enumi}{2}
\tightlist
\item
  The goods market clears. Aggregate consumption must equal aggregate
  production:
\end{enumerate}

\[
c_L+c_H=y_L+y_H.
\]

Given Walras' law, if the first two conditions are met, then the last
market automatically clears.

\subsection{g:}\label{g}

\[
\begin{aligned}
\max_{\theta,c_L,c_H,n_L,n_H}\quad
&u(c_L)-v(n_L)+u(c_H)-v(n_H)\\
\text{s.t.}\quad
&c_L=w_Ln_L-T(w_Ln_L;\theta),\\
&c_H=w_Hn_H-T(w_Hn_H;\theta),\\
&u'(c_L)w_L[1-T_y(w_Ln_L;\theta)]=v'(n_L),\\
&u'(c_H)w_H[1-T_y(w_Hn_H;\theta)]=v'(n_H),\\
&T(w_Ln_L;\theta)+T(w_Hn_H;\theta)=0,\\
&c_L+c_H=w_Ln_L+w_Hn_H.
\end{aligned}
\]

The constraints correspond to the definition of the TDCE in that they
include the two household budget constraints, the two labor supply FOCs,
the government budget balancing, and the goods market clearing.

This implies that the planner can only implement allocations under a tax
schedule from the specified family. For example, the linear Ramsey
problem imposes the same marginal rate on both the high and low types.

The Ramsey problem can be solved by solving the TDCE for a given
\(\theta\) to determine the equilibrium levels for consumption
\(c_i(\theta)\), labor supply \(n_i(\theta)\), and earnings
\(y_i(\theta)\), while balancing the government budget. Then, those
equilibrium levels should be substituted into the welfare function, to
choose the \(\theta\) that maximizes social welfare.

\breakpage

\section{Q1: Daniel}\label{q1-daniel}

\breakpage

\section{Q2 (David)}\label{q2-david}

\section{Q3 (Arun)}\label{q3-arun}

\section{Q4 (Daniel)}\label{q4-daniel}

\section{Q5 (together)}\label{q5-together}




\end{document}
